\section{Conclusion}
\label{sec:conclusion}
A rationale that cites evidence promises a good decision, true claims about the evidence, and reliance on what it cites. We audited these promises separately for LLM configuration advice, holding the evidence fixed across LLM and non-LLM advisors, ablating and deleting it, verifying explanation claims against it, and scoring decisions on independently labeled held-out outcomes. The promises came apart. Ten LLMs needed the evidence to choose well but added nothing beyond a selector that reads it; 41 of 97 audited claims misstated evidence that was validly cited, mostly alongside a correct decision; and deletion showed citation-specific reliance for only one model.

For NLP evaluation, these results suggest three practices. First, citation validity and passage-level support are not enough: an explanation can cite exactly the right items and still misreport what they contain, so attribution should be checked at the level of individual claims. Second, the errors that data-to-text research documented for tables, namely wrong numbers, superlatives, ranges and overstated regularities, persist in the free-text rationales of instruction-tuned models of every size we tested, and deserve targeted checks. Third, the value of a language model in decision support should be measured against a baseline that receives the same evidence, and decision quality and rationale correctness should be reported separately.

Two extensions would strengthen these conclusions: tasks in which the evidence is incomplete and documentation-derived priors could genuinely help, and multi-annotator or automatic verification of claims, which the structured evidence makes feasible. The protocol and released data are intended to make both measurable.
